\documentclass[12pt,a4paper]{article}
\usepackage[T1]{fontenc}
\usepackage[utf8]{inputenc}
\usepackage[french]{babel}
\usepackage{lmodern}
\usepackage{geometry}
\geometry{margin=2.5cm, top=2cm}
\usepackage{setspace}
\onehalfspacing
\usepackage{parskip}
\usepackage{enumitem}
\usepackage[hidelinks]{hyperref}
\begin{document}
\thispagestyle{empty}
\begin{minipage}[t]{0.60\textwidth}
\textbf{Cour du Québec — chambre criminelle}\\
District de Québec\\
Palais de justice de Québec\\
Salle 2.15
\end{minipage}\hfill
\begin{minipage}[t]{0.37\textwidth}
\raggedleft
Dossier : 	extbf{[N° DE DOSSIER]}\\
Accusé·e : [NOM DE L'ACCUSÉ]\\
Intimée : Sa Majesté la Procureure du Québec\\
Plaidoyer : non-culpabilité
\end{minipage}
\vspace{0.8cm}
\begin{center}
\rule{\textwidth}{0.4pt}\\[0.4em]
\textbf{\large Demande de confection d'un rapport présentenciel sous réserve de la déclaration de culpabilité et observations de la défense (article 721 (1) du Code criminel)}\\[0.3em]
Infraction reprochée : [INFRACTION REPROCHÉE]\\
Audience : 23 octobre 2026 à 9 h 30, salle 2.15, Palais de justice de Québec
\rule{\textwidth}{0.4pt}
\end{center}

\section{Objet du document}
Le document qui suit est produit en vue de l'audience du 23 octobre 2026 à 9 h 30, salle 2.15, au Palais de justice de Québec.

Son objet est double : exposer la recevabilité de la demande de confection d'un rapport présentenciel au regard de l'article 721 (1) du Code criminel et, subsidiairement, valoir observations écrites de la défense.

\textbf{Conclusion du réseau de recevabilité :} Demande de confection d'un rapport présentenciel sous réserve de la déclaration de culpabilité et observations de la défense (article 721 (1) du Code criminel) (règle appliquée : R\_NON\_COUPABLE).

\section{Faits allégués}
\subsection*{Contexte de la cause}
[À COMPLÉTER : identité de l'accusé, juridiction saisie, numéro de dossier,
infraction reprochée, date de l'incident.]

\subsection*{Faits reprochés}
[À COMPLÉTER : exposé chronologique et factuel des faits allégués par
l'accusation, tel qu'admis ou contesté par la défense.]

\subsection*{Position de la défense}
[À COMPLÉTER : plaidoyer de non-culpabilité enregistré, moyens de défense
invoqués, contestation des éléments constitutifs de l'infraction.]

\subsection*{Éléments personnels et sociaux}
[À COMPLÉTER : situation personnelle, familiale, professionnelle, état de
santé, absence ou présence d'antécédents, efforts de réparation le cas
échéant. Ces éléments seront utiles à la détermination de la peine dans
l'éventualité d'un verdict de culpabilité.]

\subsection*{Pièces et annexes}
[À COMPLÉTER : liste des pièces jointes (ex. : preuves, attestations,
correspondance), avec leur date et leur provenance.]


\section{Cadre juridique applicable}
\subsection*{Art. 721 (1) C.cr. — Rapport de l'agent de probation}
721 (1) Sous réserve des règlements d'application du paragraphe (2), lorsque l'accusé, autre qu'une organisation, plaide coupable ou est reconnu coupable d'une infraction, l'agent de probation est tenu, s'il est requis de le faire par le tribunal, de préparer et de déposer devant celui-ci un rapport écrit concernant l'accusé afin d'aider le tribunal à infliger une peine ou à décider si l'accusé devrait être absous en application de l'article 730.
\emph{Source : Code criminel (Canada), art. 721 — Rapport de l'agent de probation — https://laws-lois.justice.gc.ca/eng/acts/c-46/section-721.html (statut : hors\_ligne; activation du réseau : 1.00).}

\subsection*{Nguyen c. R., 2010 QCCA 1053}
Nguyen c. R., 2010 QCCA 1053. Demande de rapport présentenciel présentée immédiatement après le verdict : refus de première instance au motif que « le rapport présentenciel devient pas très très utile » lorsque l'accusé nie l'infraction. La Cour d'appel confirme : « rien dans la loi n'oblige le juge de première instance à demander et à obtenir un rapport présentenciel ». La demande demeure recevable après une déclaration de culpabilité, mais son accueil relève de la discrétion du tribunal.
\emph{Source : Revue de droit criminel — Nguyen c. R., 2010 QCCA 1053 — http://revuededroitcriminel.blogspot.com/2026/02/un-rapport-presentenciel-peut-etre.html (statut : hors\_ligne; activation du réseau : 1.00).}

\subsection*{Art. 721 (3) et (4) C.cr. — Contenu du rapport}
Le rapport contient, si possible : l'âge, le degré de maturité, le caractère et le comportement du délinquant ainsi que son désir de réparer le tort; ses antécédents; l'effet des mesures de rechange; tout autre renseignement exigé par le tribunal.
\emph{Source : Code criminel (Canada), art. 721 — Rapport de l'agent de probation — https://laws-lois.justice.gc.ca/eng/acts/c-46/section-721.html (statut : hors\_ligne; activation du réseau : 0.00).}

\subsection*{Art. 723 (3) C.cr. — Production d'éléments de preuve}
Le tribunal peut exiger d'office, après avoir entendu le poursuivant et le délinquant, la présentation des éléments de preuve qui pourront l'aider à déterminer la peine.
\emph{Source : Code criminel (Canada), art. 723 — Production d'éléments de preuve — https://laws-lois.justice.gc.ca/eng/acts/c-46/section-723.html (statut : hors\_ligne; activation du réseau : 0.00).}


\section{Analyse de recevabilité}
\textbf{Règle appliquée :} R\_NON\_COUPABLE.

\textbf{Motif :}

\begin{enumerate}[leftmargin=*]
\item Le plaidoyer est de non-culpabilité : la condition de l'article 721 (1) du Code criminel n'est pas réunie, aucun rapport présentenciel ne pouvant être ordonné tant que l'accusé n'est pas déclaré coupable. La demande devient recevable après le verdict de culpabilité, sous la discrétion du tribunal (Nguyen c. R., 2010 QCCA 1053). Entre-temps, la défense dépose ses observations et la preuve dont elle entend soutenir sa position.
\item Condition de recevabilité : l'accusé plaide coupable ou est reconnu coupable (art. 721 (1) C.cr.).
\item Condition procédurale : le rapport n'est dû que s'il est requis par le tribunal (art. 721 (1) C.cr.).
\item Condition d'utilité : le rapport doit pouvoir éclairer la Cour (Nguyen c. R.).
\end{enumerate}

\textbf{Chemin de propagation dans le réseau (6 nœuds actifs) :}

\texttt{d\_721\_1 -> cond\_culpabilite -> juris\_nguyen -> dec\_sous\_reserve -> dec\_imm -> cond\_requisition}


\section{Demandes}
La défense demande au tribunal :

\begin{enumerate}[leftmargin=*]
\item Prendre acte du plaidoyer de non-culpabilité enregistré.
\item Dans l'éventualité d'un verdict de culpabilité, ordonner dès lors la confection d'un rapport présentenciel conformément à l'article 721 (1) du Code criminel, le tribunal n'y étant pas tenu mais pouvant y procéder sur demande.
\item À défaut, accorder à la défense le délai de coutume pour déposer ses observations écrites en vue de la détermination de la peine.
\item Déclarer recevable le présent document, qui vaut observations écrites de la défense aux termes de l'article 723 (3) du Code criminel.
\end{enumerate}

\section{Sources et pièces consultées}
Sources téléchargées et mises en relation dans la base SQLite du pipeline :

\begin{itemize}[leftmargin=*]
\item Barreau du Québec — Site officiel — \texttt{https://www.barreau.qc.ca/} (statut : hors\_ligne; affinité max : 0.000; SHA-256 : fdef5add8011a486…)
\item Barreau du Québec — Détermination de la peine — \texttt{https://www.barreau.qc.ca/media/00idra3c/penal-peine.pdf} (statut : hors\_ligne; affinité max : 0.000; SHA-256 : eb833c98ea934033…)
\item Barreau du Québec — Procédure pénale — \texttt{https://www.barreau.qc.ca/media/osyeiodv/procedure-penale.pdf} (statut : hors\_ligne; affinité max : 0.000; SHA-256 : d965858ea3391cb2…)
\item Revue de droit criminel — Nguyen c. R., 2010 QCCA 1053 — \texttt{http://revuededroitcriminel.blogspot.com/2026/02/un-rapport-presentenciel-peut-etre.html} (statut : hors\_ligne; affinité max : 0.000; SHA-256 : 38502afc82288a07…)
\item Le carnet de droit pénal — Rapports préalables à la condamnation — \texttt{https://carnetdroitpenal.ca/fr/index.php?title=Rapports\_pr\%C3\%A9alables\_\%C3\%A0\_la\_condamnation} (statut : hors\_ligne; affinité max : 0.000; SHA-256 : aaf3ece817720544…)
\item Fondation du Barreau du Québec — Se préparer pour la cour en matière pénale (guide 2025) — \texttt{https://fondationdubarreau.qc.ca/assets/documents/Guide\_Comment-se-pr\%C3\%A9parer-pour-la-cour-en-mati\%C3\%A8re-p\%C3\%A9nale-VF-2025.pdf} (statut : hors\_ligne; affinité max : 0.000; SHA-256 : 1a531a683a74c51f…)
\item Code criminel (Canada), art. 721 — Rapport de l'agent de probation — \texttt{https://laws-lois.justice.gc.ca/eng/acts/c-46/section-721.html} (statut : hors\_ligne; affinité max : 0.000; SHA-256 : b4214fa865244497…)
\item Code criminel (Canada), art. 723 — Production d'éléments de preuve — \texttt{https://laws-lois.justice.gc.ca/eng/acts/c-46/section-723.html} (statut : hors\_ligne; affinité max : 0.000; SHA-256 : ae8827f7896b2d13…)
\item Code criminel (Légis Québec, version française) — \texttt{https://www.legisquebec.gouv.qc.ca/fr/document/lc/C-46} (statut : hors\_ligne; affinité max : 0.000; SHA-256 : e94b5dfe13448601…)
\item Règles de pratique de la cour du Québec (chambre criminelle et pénale) — \texttt{https://www.canlii.org/fr/qc/legis/regl/rlrq-1981-c-t-16-r-6/derniere/rlrq-1981-c-t-16-r-6.html} (statut : hors\_ligne; affinité max : 0.000; SHA-256 : e3b0c44298fc1c14…)
\end{itemize}

Le dépôt effectif au greffe (nombre de copies, format numérique, délai) doit être confirmé auprès du greffe de la juridiction saisie.

\section{Signature et dépôt}
\vspace{1cm}
\begin{flushright}
Fait à Québec, le 23 octobre 2026.

\vspace{1.2cm}
\rule{6cm}{0.4pt}
[NOM DE L'AVOCAT·E], membre du Barreau du Québec
Pour la défense
\end{flushright}

\vspace{0.8cm}
\hrule
\vspace{0.4em}
\small\emph{Document assemblé automatiquement par le pipeline \texttt{pipeline\_rapport\_presentenciel.py} à partir de la base SQLite \texttt{rapport\_presentenciel.db} (réseau de recevabilité, art. 721 (1) C.cr.) et de l'extrait factuel fourni. Il doit être revu par un membre du Barreau du Québec avant tout dépôt au greffe.}

\end{document}
